% =========================================================
% QUESTION 1
% =========================================================
\section{پرسش نخست}
\subsection{}
در واقع هر سیستم عامل به صورت پیش‌فرض توقع دارد که دسترسی و کنترل انحصارا کاملی به منابع سخت‌افزاری داشته باشد و برای 
اشتراک منابع طراحی نشده است، حال هم‌اکنون در حالی که هر کدام از سیستم‌های عامل مهمان توقع دسترسی
به تمامی منابع، شامل پردازنده، حافظهٔ فیزیکی و \lr{I/O} را دارد،
و همزمان درخواست‌های متعددی به این منابع مشترک می‌فرستند و این تداخل ناگهانی و زیاد
باعث تداخل فعالیت و خطا در خروجی‌های این منابع می‌باشد.
\par
برای مثال، دو سیستم عامل به طور همزمان به یک دستگاه \lr{I/O} در خواست می‌دهند و این درخواست‌ها
می‌توانند متفاوت باشند که باعث می‌شود داده‌های متناقضی بازگردانده شود، یا در صفحات حافظه ممکن است داده‌های متناقضی نوشته شود
و یا دستورات متناقضی به پردازنده ارسال گردد و این موارد می‌توانند باعث ناامنی در اجرای این چند سیستم عامل گردد.
\par
در نتیجه، این ساز و کار نیازمند یک لایهٔ واسط به نام Hypervisor می‌باشد
که منابع را به صورت زمان‌بندی شده بین این سیستم‌های عامل مهمان تقسیم می‌کند تا از دسترسی‌های تداخلی این سیستم‌های عامل جلوگیری به عمل آورد.
\par
به عبارتی، این لایهٔ واسط وظیفهٔ تخصیص کنترل‌شدهٔ منابع مختلف را تضمین می‌کند؛ صفحه‌های مختلف حافظه به سیستم‌های عامل
متفاوتی تخصیص می‌یابد، پردازنده به صورت زمان‌بندی شده به این سیستم‌های عامل تخصیص می‌یابد، دستورات \lr{I/O}
به طوری شبیه‌سازی می‌گردد تا تداخلی بین درخواست‌ها اتفاق نیفتد.

\subsection{}
اگر هستهٔ یک سیستم عامل مهمان دسترسی هستهٔ سیستم عامل میزبان را داشته باشد، آن‌گاه می‌تواند در حد دسترسی تام به منابع سخت‌افزاری 
دسترسی داشته باشد و منابع مشترک با سایر سیستم‌های عامل مهمان را به تنهایی استفاده کند و دسترسی سایرین را قطع کند.
در این حالت ماشین مجازی دیگر ایزوله نبوده و می‌تواند تنظیمات سایر سیستم‌های عامل را تغییر بدهد و منابعی که به سایر سیستم‌های عامل می‌رسد را مختل کند و مستقیما بدون لایهٔ واسطه به آنان دسترسی پیدا کند؛
این همان شرایطی می‌باشد که در طی آن یک سیستم عامل مهمان بدلیل دسترسی بیش از حد می‌تواند موجب Crash شده و اثر مخربی بر تمامی عملکرد‌های سیستم‌های 
عامل مهمان و میزبان بگذارد.

\subsection{}
در روش «تله و شبیه‌سازی»\footnote{\lr{Trap-and-Emulate}} سیستم‌عامل مهمان به طور پیش‌فرض در سطح دسترسی کاربر فعالیت می‌کند،
هنگامی که نیاز دارد دستورات با دسترسی ویژه\footnote{\lr{Privileged Instruction}} را اجرا کند،
درخواست به عنوان یک استثنا پرتاب می‌کند\footnote{\lr{Throwing Exception}} که این استثنا
توسط نرم‌افزار ناظر\footnote{Hypervisor} دریافت می‌شود و ماشین مجازی را تعلیق می‌کند
و درخواست این ماشین را به صورت شبیه‌سازی شده با دسترسی هسته اجرا می‌کند و تا هنگام نتیجه این تعلیق ادامه دارد و سپس پس از
دریافت خروجی، پاسخ درخواست را به ماشین مجازی داده و اجرای آن را از سر می‌گیرد.
\par

% =========================================================
% QUESTION 2
% =========================================================
\section{توضیح ماشین حالت فرستنده و نقش سیگنال‌های \lr{tx} و \lr{busy}}
پیمانه\LTRfootnote{Module} فرستنده بر پایه یک ماشین حالت متناهی\LTRfootnote{Finite State Machine (FSM)} با ۵ وضعیت اصلی طراحی شده است. به منظور تطبیق سرعت فرستنده و گیرنده، یک شمارنده داخلی (\lr{tick\_cnt}) در تمام وضعیت‌ها وظیفه دارد تا ماندگاری در هر بیت را به اندازه پارامتر \lr{BIT\_TICKS} تضمین نماید.

\subsection{تحلیل وضعیت‌های ماشین حالت}
گذر حالت‌ها در این پیمانه به شرح زیر است:

% =========================================================
% QUESTION 3
% =========================================================


% =========================================================
% QUESTION 4
% =========================================================

